\documentclass[12pt]{amsart}
\usepackage{amsmath,amsthm,amssymb}
\usepackage{dsfont,bm,mathtools}
\usepackage[usenames,dvipsnames,svgnames,table]{xcolor}
\usepackage[top=2.3cm,bottom=2.3cm,left=2cm,right=2cm]{geometry}
\usepackage{hyperref}
\hypersetup{hidelinks}
\allowdisplaybreaks
\emergencystretch=1.5em
\definecolor{ProposalOrange}{RGB}{160,65,0}
\definecolor{RefereePurple}{RGB}{103,34,147}
\newtheorem{thm}{Theorem}[section]
\newtheorem{lem}[thm]{Lemma}
\theoremstyle{remark}
\newtheorem{rem}[thm]{Remark}
\theoremstyle{definition}
\newtheorem{nota}[thm]{Notation}
\definecolor{ResolutionGreen}{RGB}{0,100,55}
\newcommand{\C}{\mathds{C}}
\newcommand{\Q}{\mathds{Q}}
\newcommand{\rv}{\operatorname{rv}}
\newcommand{\Gal}{\operatorname{Gal}}
\newcommand{\Aut}{\operatorname{Aut}}
\newcommand{\RV}{\mathrm{RV}}
\newcommand{\RES}{\mathrm{RES}}
\newcommand{\VF}{\mathrm{VF}}
\newcommand{\K}{\mathds{k}}
\newcommand{\Z}{\mathds{Z}}
\newcommand{\N}{\mathds{N}}
\newcommand{\vrv}{\operatorname{vrv}}
\newcommand{\GL}{\operatorname{GL}}
\newcommand{\dcl}{\operatorname{dcl}}
\newcommand{\pr}{\operatorname{pr}}
\newcommand{\defn}{\mathrm{def}}
\newcommand{\given}{\,\middle|\,}
\newcommand{\Sbase}{\mathbb S}
\newcommand{\RVlang}{\mathcal L_{\mathrm{RV}}}
\numberwithin{equation}{section}
\pagestyle{plain}

\newcommand{\A}{\mathds{A}}
\newcommand{\Gm}{\mathbb G_m}
\newcommand{\Var}{\mathrm{Var}}
\newcommand{\Kres}{K_0(\RES)}
\newcommand{\Bangres}{{!K_0(\RES)}}
\newcommand{\Kraw}{K_0^{\widehat\mu,\mathrm{raw}}(\Var_{\C})}
\newcommand{\Kmon}{K_0^{\widehat\mu}(\Var_{\C})}
\newcommand{\E}{\mathbb E_b}
\newcommand{\Vol}{\operatorname{Vol}}
\newcommand{\vv}{\operatorname{val}}
\newcommand{\LiftingTheorem}{the lifting theorem of Section~3}

\theoremstyle{plain}
\newtheorem*{thm*}{Theorem}

\newcommand{\R}{\mathds R}
\newcommand{\Sfiber}{\mathbb S_b}
\newcommand{\Rbase}[1]{{!K_0(\RES_{#1})}}
\newcommand{\IsVFVG}{(\VF,\Gamma)\text{-generated}}
\newcommand{\gsk}{K_+}
\begin{document}
\setcounter{section}{5}
\setcounter{subsection}{2}
{\small\noindent Forey--Yin, \emph{Bounded integral and motivic Milnor fiber}.
Section 5 review: R092--R093.
Approved context is black; proposed manuscript text is orange.}
\medskip

\noindent\emph{Proof of Theorem 5.17 (continued).}
The total spaces are represented by
$\bm W=(W,q)$ and $\bm W'=(W',q')$ over the mixed RV base
$S$. All base coordinates, finite tags and finite-to-one
label maps are retained.
The preceding parameter-control argument gives a finite
partition on which Gamma families depend only on base
values.

{\color{ProposalOrange}
\medskip
\noindent\emph{Support notation and further refinements.}
For a definable family of value supports put
\[
 D_{\gamma,i}=\{\delta:(\gamma,\delta)\in D_i\},\qquad
 D_i=\{(\gamma,\delta):\gamma\in E,\ \delta\in D_{\gamma,i}\}.
\]
For the primed family the corresponding notation is
$D'_{\gamma,i'}$ and $(D'_{i'})_{i'}$.
Each function $\gamma\mapsto\chi_b(D_{\gamma,i})$ has
finite range and definable level sets.
A common finite partition of $E$ makes all these functions,
and their primed counterparts, constant. On one such piece,
write the constants as $e_i,e'_{i'}$; Gamma Fubini gives
\[
 \chi_b(D_i)=e_i\chi_b(E),\qquad
 \chi_b(D'_{i'})=e'_{i'}\chi_b(E).
\]
Whenever the source or the base is restricted further,
we refine again by the level sets of the resulting fiber
Euler functions.

\medskip
\noindent\emph{A relative Gamma-push calculation.}
We use the following calculation over the original ground
in place of choosing a common fixed-value residue fiber
of the base.
Let $\bm X=(X,p)\in\RV[k]$, let
$J\subseteq X\times\Gamma^r$ be definable, and put
\[
 J_x=\{\xi:(x,\xi)\in J\},\qquad
 Y=\{(x,z):x\in X,\ \vrv(z)\in J_x\}.
\]
Give $Y$ the finite-to-one label $(p(x),z)$, of degree
$k+r$. Set $g=[\Gm]$ in $R(\Sbase)$.
The integer function $c(x)=\chi_b(J_x)$ has finite range
and definable level sets $X_e=\{x:c(x)=e\}$, and
\begin{equation}\label{relative:Gamma:push}
 \mathbb E_{\mathrm b}^{\Sbase}([Y])
       =g^r\sum_e e\,
               \mathbb E_{\mathrm b}^{\Sbase}([X_e]).
\end{equation}
The sum is finite. All classes in this formula use the
ungraded bounded retraction; the specified RV labels are
retained before that evaluation.
Two finite-to-one labels on the same RV set give the same
bounded value: pad them with constant-one coordinates to
a common degree. The underlying identity is then an RV
isomorphism. Each padding multiplies by the degree-one
point class, whose bounded value is $1$.

Here is a proof. First work on fixed zero-coordinate
strata, retaining the finite tags.
Apply the preceding parameter-control argument to the
Gamma fibers $J_x$, treating $x$ as the RV parameter.
After a finite partition of $X$, the set $J_x$ depends only
on $\vrv(x)$ through a fixed ground-definable Gamma
family. Ordinary Gamma fiber integration now proves
finite-range constructibility of $c$.
Intersect this partition with its level sets.

Apply Lemma 5.12 on each resulting part and write one chart
as
\[
 h_M(X_0)=V\times I^\sharp,\qquad
 \vrv(V)=\{\beta\},\qquad I\subseteq\Gamma^\ell.
\]
Here $\ell$ denotes the number of free RV coordinates,
rather than the dimension of the Gamma set $I$.
Because $M$ is integral unimodular, the original values
are $M^{-1}(\beta,\eta)$ when $\eta=\vrv(w)$.
The corresponding part of $Y$ therefore has the chart
\[
 V\times K^\sharp,\qquad
 K=\{(\eta,\xi):\eta\in I,\
            \xi\in J_{M^{-1}(\beta,\eta)}\}.
\]
The subscript on $J$ here denotes the value-parametrized
Gamma family on this selected part.
If $c=e$ there, Gamma Fubini gives
$\chi_b(K)=e\chi_b(I)$.
The monomial evaluation of the bounded retraction is
\[
 \mathbb E_{\mathrm b}([Y_0])
      =[V]\,g^{\ell+r}\chi_b(K)
      =g^r e\,[V]\,g^\ell\chi_b(I)
      =g^r e\,\mathbb E_{\mathrm b}([X_0]).
\]
The coordinate change on $Y_0$ is $h_M$ on $x$ and the
identity on $z$, so it preserves the finite labels by
transport. Summing the finitely many parts proves
(\ref{relative:Gamma:push}), entirely over the original
ground.

\medskip
\noindent\emph{Keep the residue families over their RV base.}
Once the finite monomial witness packages have been
descended and prepared over $S$, write the resulting
fiber charts as
\[
 h_{M_i}(U_{\tau,i})
      =V_{i,\tau}\times I_{i,\tau}^{\sharp},
 \qquad \tau\in S.
\]
The matrices, arities, finite tags, label maps and graphs
are part of these packages.
After the preceding parameter control and a finite common
partition of $S$, the numbers
$\chi_b(I_{i,\tau})=e_i$ are constant on a selected base
piece; do the same for the primed family.
The affine map
$\eta\mapsto M_i^{-1}(\beta_{i,\tau},\eta)$ identifies
this free support with the original value piece, so their
bounded Euler characteristics agree.
Retain every source chart on that piece. A chart may have
empty fibers, and its projection need not equal the
projections of other charts.

Put
\[
 \mathcal V_i=\{(\tau,v):\tau\in S,\ v\in V_{i,\tau}\}.
\]
This is an RV object with a chosen finite-to-one label;
it need not be a RES object over the ground.
Apply (\ref{relative:Gamma:push}) with
$X=\mathcal V_i$ and Gamma fiber $I_{i,\tau}$.
The maps $(\tau,u)\mapsto(\tau,h_{M_i}(u))$
retain the RV base coordinate. Thus
\begin{equation}\label{fubini:family:coefficients}
 \mathbb E_{\mathrm b}([\bm W])
   =\sum_i e_i g^{\ell_i}
               \mathbb E_{\mathrm b}([\mathcal V_i]),
 \qquad
 \mathbb E_{\mathrm b}([\bm W'])
   =\sum_{i'} e'_{i'}g^{\ell'_{i'}}
               \mathbb E_{\mathrm b}([\mathcal V'_{i'}]).
\end{equation}
The product labels in the calculation can be replaced by
the transported labels of $\bm W,\bm W'$: the monomial
evaluation in the ungraded bounded target is independent
of that label choice, as in the earlier evaluation formula.

There is a useful special case for the later RES relations.
For any definable RV family
$\mathcal Z=\{(\tau,z):z\in Z_\tau\}$ and definable
$\alpha:S\to\Gamma^n$, set
\[
 T_{\alpha,\mathcal Z}
  =\{(\tau,z,u):(\tau,z)\in\mathcal Z,\
                        \vrv(u)=\alpha(\tau)\}.
\]
Its Gamma fibers are singletons, of bounded Euler
characteristic $1$. Hence
\begin{equation}\label{fubini:relative:torsor}
 \mathbb E_{\mathrm b}([T_{\alpha,\mathcal Z}])
      =g^n\mathbb E_{\mathrm b}([\mathcal Z])
      =\mathbb E_{\mathrm b}([(\Gm)^n\times\mathcal Z]).
\end{equation}
This treats a varying value-torsor relation directly over
the RV base.

The remaining task is to obtain the finite monomial and
RES-relation witness packages over this RV base from the
hypothesis over $\Sbase\langle b\rangle$.
The calculations above do not assert that descent.
They show how to use the descended packages: combine
(\ref{fubini:family:coefficients}) with
(\ref{fubini:relative:torsor}), additivity and the actual
bijections in the finite relation witnesses.
A single finite Boolean partition of the base handles
both families and all their supports. No selected chart
is required to project surjectively, and every torus
factor in these formulas is retained.
}

\medskip
{\normalfont\small\color{RefereePurple}
\noindent\textbf{Referee R092 [Typo; revised] --
Order of the subscripts.}
Use $D_{\gamma,i}$ and $D'_{\gamma,i'}$ consistently,
including in the joint support sets, and use
$(D'_{i'})_{i'}$.\par
\medskip
\noindent\textbf{Referee R093 [Clarification; revised] --
Explain the common refinement.}
The common-refinement argument can be justified by
refining the projected supports, restricting to common
base pieces and re-establishing the fiber Euler constants.
Under the revised original-language approach, the
familywise calculation above supplies a route that
retains the residue families over the existing RV base.\par}
\medskip
{\normalfont\small\color{ProposalOrange}
\noindent\textbf{Proposed disposition of R092--R093.}
Correct the support notation and replace the common-twistback
and weighted fixed-fiber reduction before the congruence test
by the orange text. The finite witness construction and its
descent (R094--R095) remain for the next batch, adapted to
these RV families.}
\end{document}
